\chapter{pysisyphus input (menu 32)}
\label{ch:menu32}
\menulabel{32}

\section{Purpose}

pysisyphus (GPL-3.0) is an external program for exploring potential-energy
surfaces: geometry optimisations and transition-state searches with its own
optimisers (RFO/GDIIS, RS-P-RFO, dimer), chain-of-states methods and IRC
integration, driving external quantum-chemistry programs.  This menu writes
the input of the two jobs this toolkit's rescue chains use -- a minimum
optimisation and a transition-state search with ORCA as the calculator -- so
that a user can hand the geometry work to pysisyphus and read the run back in
menu 33.  Nothing of pysisyphus is linked in: the program stays external, and
only its input format (a YAML file) is written here.

\section{What you need}

A structure (an XYZ file, an ORCA input with an inline coordinate block, or an
ORCA output whose final geometry is taken), the ORCA keywords to run (any
simple line, e.g.\ \code{PBE0 D4 def2-SVP}), the charge and multiplicity, and
the machine settings (cores per call and memory per core).  A block string
(e.g.\ \code{\%scf maxiter 300 end}) can be attached for every call.

\section{Walkthrough}

The first session of the example writes a minimum optimisation with pure
defaults; the second writes a TS search (keywords and job kind given, model
Hessian by default); the third deliberately asks for the quantum-Hessian
route and is refused.

\lstinputlisting[style=fbkcode, firstline=38, lastline=62,
  title={The menu-32 example (the minimum optimisation: the record, the two
  written files and the run guidance)}]
  {examples/expected/m32_pysisyphus.txt}

The TS input takes the same shape with a \code{tsopt} block
(\code{type: rsprfo}, \code{hessian\_init: fischer}, optionally an
\code{rx\_modes} flow sequence selecting the mode to follow).  Asking for the
quantum Hessian is refused with the measured reason:

\lstinputlisting[style=fbkcode, firstline=163, lastline=178,
  title={The menu-32 example (the quantum-Hessian refusal)}]
  {examples/expected/m32_pysisyphus.txt}

\section{What you get}

\file{<stem>.pysisyphus.xyz} (the structure, standalone), \file{<stem>.pysisyphus.yaml}
(the input) and a report \file{<stem>.pysisyphus.fbk.md}.  The YAML follows the
schema pysisyphus's own example set uses and this project re-measured on
pysisyphus 1.0.0 with ORCA 6.1.1: \code{geom} (type + the XYZ file),
\code{calc} (\code{orca5}, the keywords string as the \code{!} line of every
call, charge, mult, \code{pal}, \code{mem} in MB, optional \code{blocks}), and
\code{opt} or \code{tsopt}.

\section{The quantum-Hessian boundary (measured)}

\begin{readbox}
\begin{itemize}
  \item ORCA 6 writes an extra \code{\$multiplicity} block into its
    \file{.hess} file.  The parser of the pysisyphus ORCA interface (measured
    identical in release 1.0.0 and on the master branch) has no grammar for
    that block and stops: \code{ParseException: Expected W:(0-9), found '\$'},
    followed by \code{RunAfterCalculationFailedException}.  Consequences:
    \code{hessian\_init: calc} (the \code{tsopt} default) and
    \code{do\_hess: true} are unusable with ORCA 6.1.1;
  \item everything gradient-driven is unaffected -- ORCA's own \file{.engrad}
    sidecar parses fine -- so minimum optimisations run end to end, and TS
    searches are reliable with a \emph{model} Hessian
    (\code{hessian\_init: fischer}, the menu's default; \code{lindh},
    \code{simple}, \code{swart}, \code{unit} are the other model choices);
  \item the generator therefore refuses \code{calc} by name and does not offer
    \code{do\_hess}; the frequency verification of a result stays outside
    pysisyphus -- a plain ORCA \code{Freq} job on the closing geometry
    (menu 30 is the toolkit's own route back from an imaginary mode).
\end{itemize}
\end{readbox}

\section{Running it}

\begin{itemize}
  \item the engine command is read from \file{\textasciitilde/.pysisyphusrc}
    (INI: \code{[orca5] cmd=/path/to/orca}), or from \code{\$PATH} when that
    section is absent;
  \item pysisyphus stages every calculator call in a temporary directory under
    \code{\$TMPDIR} and copies the kept files into \file{qm\_calcs/} when the
    call finishes -- point \code{\$TMPDIR} at a real scratch disk;
  \item keep the console: \code{pysis <file> > run.out}; that capture is the
    run's record (the cycle table, the outcome marker, the closing summary)
    and it is what menu 33 reads back.
\end{itemize}
